\documentclass[12pt,a4paper]{article}

% --- XeLaTeX + polskie znaki przez fontspec ---
\usepackage{fontspec}
\setmainfont{DejaVu Serif}
\setmonofont{DejaVu Sans Mono}
\usepackage{polyglossia}
\setdefaultlanguage{polish}

\usepackage{dirtree}
\usepackage{geometry}
\usepackage{hyperref}
\usepackage{listings}
\usepackage{xcolor}
\usepackage{booktabs}
\usepackage{longtable}
\usepackage{array}
\usepackage{enumitem}
\usepackage{titlesec}
\usepackage{fancyhdr}
\usepackage{parskip}
\usepackage{microtype}

\geometry{a4paper, top=2.5cm, bottom=2.5cm, left=2.5cm, right=2.5cm}

% --- Nagłówki i stopki ---
\pagestyle{fancy}
\fancyhf{}
\fancyhead[L]{\small Digital Rehab App --- Dokumentacja Techniczna}
\fancyhead[R]{\small v1.0.0}
\fancyfoot[C]{\thepage}
\renewcommand{\headrulewidth}{0.4pt}
\setlength{\headheight}{14.5pt}

% --- Styl kodu ---
\lstset{
  basicstyle=\ttfamily\small,
  backgroundcolor=\color{gray!10},
  frame=single,
  breaklines=true,
  keywordstyle=\color{blue!70},
  commentstyle=\color{green!50!black},
  stringstyle=\color{red!60},
  showstringspaces=false,
  numbers=left,
  numberstyle=\tiny\color{gray},
  tabsize=2,
  literate=
    {ą}{{\k{a}}}1 {Ą}{{\k{A}}}1
    {ę}{{\k{e}}}1 {Ę}{{\k{E}}}1
    {ó}{{\'o}}1  {Ó}{{\'O}}1
    {ś}{{\'s}}1  {Ś}{{\'S}}1
    {ł}{{\l{}}}1 {Ł}{{\L{}}}1
    {ż}{{\.z}}1  {Ż}{{\.Z}}1
    {ź}{{\'z}}1  {Ź}{{\'Z}}1
    {ć}{{\'c}}1  {Ć}{{\'C}}1
    {ń}{{\'n}}1  {Ń}{{\'N}}1
}

% --- Linki ---
\hypersetup{colorlinks=true, linkcolor=blue!60!black, urlcolor=blue!60!black}

\titlespacing*{\section}{0pt}{1.5em}{0.8em}
\titlespacing*{\subsection}{0pt}{1.2em}{0.5em}
\titlespacing*{\subsubsection}{0pt}{1em}{0.4em}

% ============================================================
\begin{document}
% --- Strona tytułowa ---
\begin{titlepage}
  \centering
  \vspace*{2cm}
  {\Huge\bfseries Digital Rehab App\\[0.4em]}
  {\Large Dokumentacja Techniczna Systemu}
  \vspace{1.5cm}

  \rule{\linewidth}{0.5pt}
  \vspace{0.8cm}

  {\large\bfseries Autorzy}\\[0.6em]
  \begin{tabular}{c}
    Bartek Czyż,
    Piotr Kotuła, 
    Michał Patronik,
    Miłosz Śmieja 
  \end{tabular}

  \vspace{0.8cm}
  \rule{\linewidth}{0.5pt}
  \vspace{0.5cm}

  \begin{tabular}{ll}
    \textbf{Wersja dokumentu:}   & 1.0.0 \\[4pt]
    \textbf{Wersja aplikacji:}   & 1.0.0+1 \\[4pt]
    \textbf{Data wydania:}       & czerwiec 2026 \\[4pt]
    \textbf{Język dokumentacji:} & polski \\[4pt]
    \textbf{Forma:}              & elektroniczna (PDF) \\
  \end{tabular}

  \vspace{1cm}
  \rule{\linewidth}{0.5pt}
  \vspace{1cm}

  {\large\bfseries Projekt Zespołowy --- PWr}\\[0.6em]
  {\large Politechnika Wrocławska}\\
  {\large Wydział Informatyki i Telekomunikacji}

  \vspace{2cm}
  \begin{tabular}{ll}
    \textbf{Typ dokumentu:} & Dokumentacja powykonawcza \\[4pt]
    \textbf{Status:}        & Wersja deweloperska \\
  \end{tabular}

  \vfill
  {\small Dokumentacja sporządzona zgodnie z wymaganiami dot.\ dokumentacji\\
  systemów informatycznych (Zał.\ nr 1.1 do SIWZ oraz Zał.\ nr 6 do Umowy)}
\end{titlepage}

\tableofcontents
\newpage

% ============================================================
\section{Informacje ogólne}

\subsection{Cel dokumentu}

Niniejszy dokument stanowi kompletną dokumentację techniczną systemu \textbf{Digital Rehab App} --- aplikacji mobilnej wspomagającej rehabilitację i ćwiczenia fizyczne. Dokumentacja obejmuje opis architektury systemu, instrukcje instalacji i konfiguracji, opis funkcjonalności, procedury eksploatacyjne oraz informacje o licencjonowaniu komponentów.

Dokument przeznaczony jest dla:
\begin{itemize}
  \item administratorów systemu odpowiedzialnych za instalację i utrzymanie,
  \item użytkowników końcowych (klientów i trenerów/fizjoterapeutów),
  \item programistów rozwijających lub integrujących system z innymi rozwiązaniami.
\end{itemize}

\subsection{Zakres systemu}

Digital Rehab App jest systemem dwuwarstwowym składającym się z:
\begin{itemize}
  \item \textbf{Aplikacji mobilnej} (frontend) --- zbudowanej w technologii Flutter, uruchamianej na urządzeniach z systemem Android (API~24+),
  \item \textbf{Serwera REST API} (backend) --- zbudowanego w technologii Spring Boot, komunikującego się z relacyjną bazą danych MySQL.
\end{itemize}

System umożliwia pacjentom samodzielne wykonywanie ćwiczeń rehabilitacyjnych z detekcją ruchu opartą na uczeniu maszynowym (Google ML~Kit), a trenerom/fizjoterapeutom --- monitorowanie postępów swoich podopiecznych.

\subsection{Wersjonowanie dokumentacji}

\begin{center}
\begin{tabular}{cccp{5.5cm}}
  \toprule
  \textbf{Nr wersji} & \textbf{Data} & \textbf{Autor} & \textbf{Opis zmian} \\
  \midrule
  1.0.0 & 2026-06 & Zespół projektowy & Pierwsza wersja dokumentacji końcowej \\
  \bottomrule
\end{tabular}
\end{center}

% ============================================================
\section{Architektura systemu}

\subsection{Architektura ogólna}

System Digital Rehab App oparty jest na architekturze \textbf{klient--serwer} z podziałem na trzy warstwy:
\begin{enumerate}
  \item \textbf{Warstwa prezentacji} --- aplikacja Flutter na urządzeniu mobilnym użytkownika,
  \item \textbf{Warstwa logiki biznesowej} --- serwer Spring Boot eksponujący REST API,
  \item \textbf{Warstwa danych} --- relacyjna baza danych MySQL.
\end{enumerate}

Komunikacja między klientem a serwerem odbywa się poprzez protokół HTTP z wymianą danych w formacie JSON. Dane dotyczące historii treningów przechowywane są lokalnie na urządzeniu mobilnym w bazie Hive (NoSQL, offline-first) --- synchronizacja z backendem jest planowana w kolejnej wersji systemu.

\subsection{Schemat blokowy systemu}

\begin{center}
\renewcommand{\arraystretch}{1.4}
\begin{tabular}{|c|}
\hline
\textbf{Urządzenie mobilne (Android)} \\
Flutter 3.41+ / Dart 3.11+ \\
\hline
$\downarrow\uparrow$ \ HTTP REST (JSON) \ $\downarrow\uparrow$ \\
\hline
\textbf{Serwer aplikacji} \\
Spring Boot 4.0.6 / Java 21 \ -- \ \texttt{localhost:8080} \\
\hline
$\downarrow\uparrow$ \ JDBC / JPA (Hibernate) \ $\downarrow\uparrow$ \\
\hline
\textbf{Baza danych MySQL 8} \\
\texttt{localhost:3306 / rehab\_db} \\
\hline
\end{tabular}
\end{center}

\subsection{Komponenty frontendu (Flutter)}

\begin{center}
\begin{tabular}{lp{8cm}}
  \toprule
  \textbf{Komponent} & \textbf{Opis} \\
  \midrule
  \texttt{lib/config/}        & Konfiguracja URL backendu i kluczy API \\
  \texttt{lib/core/}          & Motyw graficzny aplikacji \\
  \texttt{lib/data/models/}   & Modele danych (Exercise, WorkoutSession) \\
  \texttt{lib/data/services/} & Warstwa serwisów REST (ApiService) \\
  \texttt{lib/ui/screens/}    & Ekrany panelu klienta i trenera \\
  \texttt{lib/ui/widgets/}    & Wielokrotnie używane komponenty UI \\
  \texttt{lib/vision/}        & Silnik detekcji pozy (ML Kit + licznik powtórzeń) \\
  \bottomrule
\end{tabular}
\end{center}

\subsection{Komponenty backendu (Spring Boot)}

\begin{center}
\begin{tabular}{lp{8cm}}
  \toprule
  \textbf{Komponent} & \textbf{Opis} \\
  \midrule
  \texttt{controllers/}    & Warstwa REST: UsersController, ExercisesController \\
  \texttt{entities/}       & Encje JPA: Users, Exercises, ExercisesHistory \\
  \texttt{repositories/}   & Interfejsy Spring Data JPA \\
  \texttt{DataInitializer} & Inicjalizacja bazy przykładowymi danymi \\
  \texttt{application.properties} & Konfiguracja połączenia z bazą danych \\
  \bottomrule
\end{tabular}
\end{center}

% ============================================================
\section{Wymagania systemowe}

\subsection{Wymagania serwera (backend)}

\begin{center}
\begin{tabular}{lll}
  \toprule
  \textbf{Komponent} & \textbf{Wersja minimalna} & \textbf{Rekomendowana} \\
  \midrule
  Java (JDK)  & 17     & 21 (LTS) \\
  Maven       & 3.6    & 3.9+ \\
  MySQL       & 8.0    & 8.0 \\
  Docker      & dowolna aktualna & dowolna aktualna \\
  RAM         & 512 MB & 1 GB \\
  Dysk        & 500 MB & 1 GB \\
  \bottomrule
\end{tabular}
\end{center}

\subsection{Wymagania urządzenia mobilnego (frontend)}

\begin{center}
\begin{tabular}{lll}
  \toprule
  \textbf{Komponent} & \textbf{Wymaganie} & \textbf{Uwagi} \\
  \midrule
  System operacyjny & Android API 24 (7.0+) & Wymagane do ML Kit \\
  Flutter SDK       & 3.41+                 & Do kompilacji \\
  Dart SDK          & 3.11+                 & Do kompilacji \\
  Kamera            & Wymagana              & Detekcja pozy \\
  RAM               & min. 2 GB             & 3 GB rekomendowane \\
  \bottomrule
\end{tabular}
\end{center}

% ============================================================
\newpage
\section{Instalacja i konfiguracja}

\subsection{Baza danych MySQL}

Zalecanym sposobem uruchomienia bazy danych jest kontener Docker. Poniższe polecenie tworzy instancję MySQL~8:

\begin{lstlisting}[language=bash, caption={Pierwsze uruchomienie bazy danych (Docker)}]
docker run -d \
  --name rehab_mysql \
  -e MYSQL_ALLOW_EMPTY_PASSWORD=yes \
  -e MYSQL_DATABASE=rehab_db \
  -p 3306:3306 \
  mysql:8
\end{lstlisting}

Przy kolejnych uruchomieniach (kontener już istnieje):
\begin{lstlisting}[language=bash, caption={Ponowne uruchomienie kontenera}]
docker start rehab_mysql
\end{lstlisting}

Alternatywnie, przy lokalnej instalacji MySQL bez Dockera:
\begin{lstlisting}[language=bash, caption={Konfiguracja bez Dockera}]
sudo systemctl start mysql
mysql -u root -e "CREATE DATABASE IF NOT EXISTS rehab_db;"
\end{lstlisting}

\subsection{Konfiguracja backendu}

Parametry połączenia z bazą danych znajdują się w pliku \texttt{backend/src/main/resources/application.properties}:

\begin{lstlisting}[caption={Domyślna konfiguracja application.properties}]
spring.application.name=aplikacja_wspomagajaca_cwiczenia
spring.datasource.url=jdbc:mysql://localhost:3306/rehab_db
spring.datasource.username=root
spring.datasource.password=
spring.jpa.hibernate.ddl-auto=update
spring.jpa.properties.hibernate.dialect=org.hibernate.dialect.MySQLDialect
\end{lstlisting}

Przy wdrożeniu na innym hoście należy zmienić wartość \texttt{spring.datasource.url} oraz uzupełnić dane uwierzytelniające.

\subsection{Uruchomienie backendu}

\begin{lstlisting}[language=bash, caption={Uruchomienie serwera Spring Boot}]
cd backend
chmod +x mvnw
./mvnw spring-boot:run
\end{lstlisting}

Serwer startuje na porcie \texttt{8080}. Przy pierwszym uruchomieniu \texttt{DataInitializer} automatycznie tworzy schemat bazy (Hibernate DDL) i wypełnia ją danymi demonstracyjnymi.

\subsection{Konfiguracja i uruchomienie frontendu}

Przed kompilacją należy skopiować i uzupełnić plik konfiguracyjny:
\begin{lstlisting}[language=bash, caption={Konfiguracja pliku app\_config.dart}]
cp frontend/lib/config/app_config.example.dart \
   frontend/lib/config/app_config.dart
\end{lstlisting}

Następnie w \texttt{app\_config.dart} ustawić wartość \texttt{backendUrl}:

\begin{center}
\begin{tabular}{ll}
  \toprule
  \textbf{Scenariusz} & \textbf{Wartość backendUrl} \\
  \midrule
  Emulator Android   & \texttt{http://10.0.2.2:8080} \\
  Fizyczne urządzenie & \texttt{http://<IP\_MASZYNY>:8080} \\
  \bottomrule
\end{tabular}
\end{center}

IP maszyny: \texttt{hostname -I | awk '\{print \$1\}'}.

Dodatkowo należy uzupełnić \texttt{rapidApiKey} --- klucz generowany na \href{https://rapidapi.com/justin-tooke/api/exercisedb}{RapidAPI (ExerciseDB)}, wymagany do katalogu ćwiczeń.

\begin{lstlisting}[language=bash, caption={Kompilacja i uruchomienie frontendu}]
cd frontend
flutter pub get
flutter run
# lub na konkretnym urządzeniu:
flutter devices
flutter run -d <ID_URZĄDZENIA>
\end{lstlisting}

% ============================================================
\section{Schemat bazy danych}

\subsection{Tabela \texttt{users}}

\begin{center}
\begin{tabular}{llll}
  \toprule
  \textbf{Kolumna} & \textbf{Typ} & \textbf{Ograniczenia} & \textbf{Opis} \\
  \midrule
  \texttt{id}          & BIGINT  & PK, AUTO\_INCREMENT     & Identyfikator użytkownika \\
  \texttt{email}       & VARCHAR & UNIQUE, NOT NULL        & Adres e-mail (login) \\
  \texttt{password}    & VARCHAR & NOT NULL                & Hasło (patrz sekcja \ref{sec:ograniczenia}) \\
  \texttt{role}        & VARCHAR & NOT NULL                & \texttt{KLIENT} lub \texttt{TRENER} \\
  \texttt{trainer\_id} & BIGINT  & FK $\rightarrow$ users(id) & ID trenera (nullable) \\
  \bottomrule
\end{tabular}
\end{center}
\newpage
\subsection{Tabela \texttt{exercises}}

\begin{center}
\begin{tabular}{llll}
  \toprule
  \textbf{Kolumna} & \textbf{Typ} & \textbf{Ograniczenia} & \textbf{Opis} \\
  \midrule
  \texttt{id}          & BIGINT  & PK, AUTO\_INCREMENT & Identyfikator ćwiczenia \\
  \texttt{name}        & VARCHAR & NOT NULL            & Nazwa ćwiczenia \\
  \texttt{description} & VARCHAR & --                  & Opis i wskazówki wykonania \\
  \bottomrule
\end{tabular}
\end{center}

\subsection{Tabela \texttt{exercise\_history}}

\begin{center}
\begin{tabular}{llll}
  \toprule
  \textbf{Kolumna} & \textbf{Typ} & \textbf{Ograniczenia} & \textbf{Opis} \\
  \midrule
  \texttt{id}           & BIGINT   & PK, AUTO\_INCREMENT        & Identyfikator wpisu \\
  \texttt{user\_id}     & BIGINT   & FK $\rightarrow$ users(id) & Kto ćwiczył \\
  \texttt{exercise\_id} & BIGINT   & FK $\rightarrow$ exercises(id) & Jakie ćwiczenie \\
  \texttt{repetitions}  & INT      & NOT NULL                   & Liczba powtórzeń \\
  \texttt{accuracy}     & INT      & 0--100                     & Poprawność (\%) \\
  \texttt{created\_at}  & DATETIME & DEFAULT NOW()              & Data i czas sesji \\
  \bottomrule
\end{tabular}
\end{center}

\subsection{Relacje między tabelami}

\begin{itemize}
  \item \texttt{users.trainer\_id} $\rightarrow$ \texttt{users.id} --- relacja wiele-do-jednego (klienci do trenera),
  \item \texttt{exercise\_history.user\_id} $\rightarrow$ \texttt{users.id} --- historia przypisana do użytkownika,
  \item \texttt{exercise\_history.exercise\_id} $\rightarrow$ \texttt{exercises.id} --- historia przypisana do ćwiczenia.
\end{itemize}

% ============================================================
\section{Dokumentacja API}

\subsection{Informacje ogólne}

\begin{center}
\begin{tabular}{ll}
  \toprule
  \textbf{Parametr} & \textbf{Wartość} \\
  \midrule
  Protokół      & HTTP \\
  Format danych & JSON \\
  Port          & 8080 \\
  Base URL      & \texttt{http://<HOST>:8080} \\
  Autentykacja  & Brak (TODO: JWT) \\
  \bottomrule
\end{tabular}
\end{center}

\newpage
\subsection{Endpointy --- Użytkownicy}

\subsubsection{Rejestracja nowego konta}

\begin{lstlisting}[caption={POST /api/users/newAccount}]
POST /api/users/newAccount
Content-Type: application/json

{
  "email": "jan@example.com",
  "password": "haslo123",
  "role": "KLIENT",
  "trainer": { "id": 1 }   // opcjonalne - ID trenera
}
\end{lstlisting}

Odpowiedzi: \textbf{200~OK} (sukces, zwraca obiekt użytkownika), \textbf{400~Bad~Request} (niepoprawny email/hasło/rola lub nieistniejący trener), \textbf{409~Conflict} (email już zajęty).

\subsubsection{Logowanie}

\begin{lstlisting}[caption={POST /api/users/login}]
POST /api/users/login
Content-Type: application/json

{ "email": "jan@example.com", "password": "haslo123" }
\end{lstlisting}

Odpowiedzi: \textbf{201~Created} (obiekt użytkownika z polem \texttt{role}), \textbf{400~Bad~Request} (brak emaila), \textbf{403~Forbidden} (błędne hasło).

\subsubsection{Lista pacjentów trenera}
\begin{lstlisting}[caption={GET /trainers/\{id\}/clients}]
GET /trainers/{id}/clients
\end{lstlisting}
Zwraca listę użytkowników o roli \texttt{KLIENT} przypisanych do trenera.

\subsubsection{Postępy pacjenta}
\begin{lstlisting}[caption={GET /users/\{id\}/progress}]
GET /users/{id}/progress
\end{lstlisting}

\subsection{Endpointy --- Ćwiczenia}

\subsubsection{Lista ćwiczeń}
\begin{lstlisting}[caption={GET /api/exercises}]
GET /api/exercises
\end{lstlisting}

\subsubsection{Zapisanie wyniku treningu}
\begin{lstlisting}[caption={POST /api/history}]
POST /api/history
Content-Type: application/json

{
  "user":     { "id": 2 },
  "exercise": { "id": 1 },
  "repetitions": 15,
  "accuracy":    87
}
\end{lstlisting}

\subsubsection{Historia ćwiczeń użytkownika}
\begin{lstlisting}[caption={GET /users/\{id\}/history}]
GET /users/{id}/history
\end{lstlisting}

% ============================================================
\section{Dokumentacja użytkownika}

\subsection{Uruchomienie i logowanie}

Po uruchomieniu aplikacji wyświetlany jest ekran powitalny (\texttt{WelcomeScreen}), a następnie ekran logowania. Użytkownik podaje adres e-mail, hasło oraz wybiera rolę (Klient / Trener/Fizjoterapeuta). Po poprawnym zalogowaniu aplikacja przekierowuje do odpowiedniego panelu.

\textbf{Dane demonstracyjne dostępne po instalacji:}

\begin{center}
\begin{tabular}{lll}
  \toprule
  \textbf{Rola} & \textbf{E-mail} & \textbf{Hasło} \\
  \midrule
  Trener        & \texttt{trener@test.com}            & \texttt{trener123} \\
  Klient        & \texttt{klient@test.com}            & \texttt{klient123} \\
  Pacjent demo  & \texttt{marek.kowalski@demo.pl}     & \texttt{demo1234} \\
  Pacjent demo  & \texttt{anna.nowak@demo.pl}         & \texttt{demo1234} \\
  Pacjent demo  & \texttt{piotr.wisniewski@demo.pl}   & \texttt{demo1234} \\
  Pacjent demo  & \texttt{katarzyna.wojcik@demo.pl}   & \texttt{demo1234} \\
  Pacjent demo  & \texttt{tomasz.lewandowski@demo.pl} & \texttt{demo1234} \\
  \bottomrule
\end{tabular}
\end{center}

\subsection{Panel klienta}

\begin{itemize}
  \item \textbf{Katalog ćwiczeń} (\texttt{CatalogScreen}) --- przeglądanie listy ćwiczeń z zewnętrznego API ExerciseDB oraz bazy lokalnej, filtrowanie po partii mięśniowej.
  \item \textbf{Szczegóły ćwiczenia} (\texttt{ExerciseDetailScreen}) --- opis, wskazówki techniczne, przycisk rozpoczęcia sesji z kamerą.
  \item \textbf{Sesja treningowa z kamerą} (\texttt{CameraScreen}) --- kamera analizuje ułożenie ciała w czasie rzeczywistym przy użyciu Google ML~Kit Pose Detection. Silnik \texttt{PushUpCounter} zlicza powtórzenia na podstawie kąta stawu łokciowego (próg fazy dolnej:~110°, górnej:~160°).
  \item \textbf{Podsumowanie sesji} (\texttt{SessionSummaryScreen}) --- liczba powtórzeń i wskaźnik poprawności po zakończeniu treningu.
  \item \textbf{Historia treningów} (\texttt{WorkoutHistoryScreen}) --- lista minionych sesji z wykresami (fl\_chart). Dane przechowywane lokalnie (Hive).
\end{itemize}

\subsection{Panel trenera / fizjoterapeuty}

\begin{itemize}
  \item \textbf{Lista pacjentów} (\texttt{PatientListScreen}) --- przegląd podopiecznych ze statusem: aktywny / nieobecny / zgłoszony ból.
  \item \textbf{Szczegóły pacjenta} (\texttt{PatientDetailScreen}) --- telemetria, plan ćwiczeń, historia sesji pacjenta.
  \item \textbf{Edytor planu ćwiczeń} (\texttt{PlanEditorScreen}) --- tworzenie i modyfikowanie indywidualnych planów rehabilitacyjnych.
  \item \textbf{Przegląd wideo} (\texttt{VideoReviewScreen}) --- \textit{w budowie} (planowane w kolejnej wersji).
  \item \textbf{Wiadomości} --- \textit{w budowie} (planowane w kolejnej wersji).
\end{itemize}

% ============================================================
\section{Licencjonowanie}

\subsection{Oprogramowanie własne}

Kod źródłowy aplikacji (frontend Flutter i backend Spring Boot) jest własnością zespołu projektowego. Nie jest objęty żadną publiczną licencją open-source --- wszelkie prawa zastrzeżone w zakresie projektu akademickiego.

\subsection{Zależności zewnętrzne --- Backend}

\begin{center}
\begin{tabular}{lll}
  \toprule
  \textbf{Biblioteka} & \textbf{Wersja} & \textbf{Licencja} \\
  \midrule
  Spring Boot               & 4.0.6  & Apache 2.0 \\
  Hibernate ORM             & 6.x    & LGPL 2.1 \\
  MySQL Connector/J         & 8.x    & GPL 2.0 (FOSS exception) \\
  Lombok                    & 1.18.x & MIT \\
  Apache Commons Validator  & 1.10.1 & Apache 2.0 \\
  \bottomrule
\end{tabular}
\end{center}

\subsection{Zależności zewnętrzne --- Frontend}

\begin{center}
\begin{tabular}{lll}
  \toprule
  \textbf{Pakiet} & \textbf{Wersja} & \textbf{Licencja} \\
  \midrule
  Flutter SDK                    & 3.41+  & BSD 3-Clause \\
  google\_mlkit\_pose\_detection & 0.14.0 & BSD / Google \\
  camera                         & 0.11.0 & BSD 3-Clause \\
  hive / hive\_flutter           & 2.2.3  & Apache 2.0 \\
  fl\_chart                      & 0.69.0 & MIT \\
  http                           & 1.6.0  & BSD 3-Clause \\
  google\_fonts                  & 6.2.1  & Apache 2.0 \\
  permission\_handler            & 11.3.1 & MIT \\
  \bottomrule
\end{tabular}
\end{center}

\subsection{Zewnętrzne usługi / API}

ExerciseDB (RapidAPI) --- zewnętrzne API katalogu ćwiczeń. Wymaga klucza API (\texttt{rapidApiKey}). Licencja komercyjna z planem darmowym (limit zapytań miesięcznie).

% ============================================================
\section{Procedury eksploatacyjne}

\subsection{Kolejność uruchamiania komponentów}

Kolejność jest istotna:
\begin{enumerate}
  \item Uruchomić kontener MySQL: \texttt{docker start rehab\_mysql}
  \item Uruchomić backend: \texttt{./mvnw spring-boot:run} (katalog \texttt{backend/})
  \item Uruchomić aplikację mobilną: \texttt{flutter run} (katalog \texttt{frontend/})
\end{enumerate}

\subsection{Zatrzymanie systemu}

\begin{enumerate}
  \item Zamknąć aplikację mobilną
  \item Zatrzymać backend (\texttt{Ctrl+C})
  \item Zatrzymać MySQL: \texttt{docker stop rehab\_mysql}
\end{enumerate}

\subsection{Procedura backup i odtworzenia bazy danych}

Wykonanie kopii zapasowej:
\begin{lstlisting}[language=bash, caption={Kopia zapasowa bazy danych}]
docker exec rehab_mysql mysqldump -u root rehab_db \
  > backup_$(date +%Y%m%d).sql
\end{lstlisting}

Odtworzenie z kopii:
\begin{lstlisting}[language=bash, caption={Odtworzenie bazy z kopii zapasowej}]
docker exec -i rehab_mysql mysql -u root rehab_db \
  < backup_20260601.sql
\end{lstlisting}

Dane historii treningów przechowywane lokalnie na urządzeniu (Hive) można wyeksportować przez standardowy mechanizm kopii systemu Android (ADB backup).

\subsection{Weryfikacja działania systemu}

\begin{lstlisting}[language=bash, caption={Weryfikacja działania API}]
# Pobranie listy cwiczen -- oczekiwana odpowiedz: tablica JSON
curl http://localhost:8080/api/exercises
\end{lstlisting}

\subsection{Logi aplikacji}

\begin{lstlisting}[language=bash, caption={Uruchomienie backendu z zapisem logów}]
mkdir -p logs
./mvnw spring-boot:run > logs/backend.log 2>&1 &
\end{lstlisting}

Logi Hibernate (DDL, SQL) dostępne na standardowym wyjściu konsoli serwera.

\subsection{Resetowanie bazy danych do stanu początkowego}

\begin{lstlisting}[language=bash, caption={Reset bazy danych}]
docker stop rehab_mysql && docker rm rehab_mysql
docker run -d --name rehab_mysql \
  -e MYSQL_ALLOW_EMPTY_PASSWORD=yes \
  -e MYSQL_DATABASE=rehab_db \
  -p 3306:3306 mysql:8
# DataInitializer zaladuje dane po uruchomieniu backendu
./mvnw spring-boot:run
\end{lstlisting}

% ============================================================
\section{Testy jednostkowe backendu}

\subsection{Strategia testowania}

Warstwa testów jednostkowych backendu opiera się na izolowanym testowaniu kontrolerów REST
przy użyciu adnotacji \texttt{@WebMvcTest}. Podejście to uruchamia wyłącznie warstwę sieciową
(serwlet, routowanie, serializację JSON) bez inicjalizacji bazy danych --- repozytoria
zastępowane są atrapami za pomocą adnotacji \texttt{@MockitoBean}.

Testy pisane są według stylu BDD (\textit{Behavior-Driven Development}): warunek wstępny
konfigurowany przez \texttt{BDDMockito.given(\ldots).willReturn(\ldots)}, weryfikacja
odpowiedzi przez \texttt{MockMvc} i asercje \texttt{jsonPath}.

\subsection{Technologie i adnotacje}

\begin{center}
\begin{tabular}{lp{7.5cm}}
  \toprule
  \textbf{Technologia / adnotacja} & \textbf{Rola w testach} \\
  \midrule
  \texttt{@WebMvcTest(XyzController.class)} & Ładuje wyłącznie warstwę MVC wybranego kontrolera \\
  \texttt{@MockitoBean}           & Rejestruje atrapę repozytorium w kontekście Springa \\
  \texttt{MockMvc}                & Symuluje żądania HTTP bez uruchamiania serwera \\
  \texttt{BDDMockito.given()}     & Konfiguruje zachowanie atrapy (styl BDD) \\
  \texttt{jsonPath(\$.pole)}      & Weryfikuje pola w odpowiedzi JSON \\
  JUnit~5 + \texttt{@DisplayName} & Czytelne nazwy testów po polsku \\
  \bottomrule
\end{tabular}
\end{center}

\textbf{Uwaga Spring Boot 4.x:} Adnotacja \texttt{@MockBean} (Spring Boot 3.x) została
zastąpiona przez \texttt{@MockitoBean}
(\texttt{org.springframework.test.context.bean.override.mockito}).
Analogicznie \texttt{@WebMvcTest} pochodzi z nowego pakietu
\texttt{org.springframework.boot.webmvc.test.autoconfigure}.

\subsection{Uruchamianie testów}

\begin{lstlisting}[language=bash, caption={Uruchomienie wszystkich testów jednostkowych}]
cd backend
./mvnw test
\end{lstlisting}

\begin{lstlisting}[language=bash, caption={Uruchomienie testów wybranego kontrolera}]
./mvnw test -Dtest=UserControllerTest
./mvnw test -Dtest=ExercisesControllerTest
\end{lstlisting}

Raporty XML i TXT zapisywane są do katalogu \texttt{backend/target/surefire-reports/}.

\subsection{Pokryte scenariusze --- \texttt{UserControllerTest}}

Klasa testuje endpointy \texttt{POST /api/users/newAccount},
\texttt{POST /api/users/login} oraz \texttt{GET /trainers/\{id\}/clients}.

\begin{center}
\begin{longtable}{p{7.5cm}cc}
  \toprule
  \textbf{Scenariusz} & \textbf{Endpoint} & \textbf{Status} \\
  \midrule
  \endfirsthead
  \toprule
  \textbf{Scenariusz} & \textbf{Endpoint} & \textbf{Status} \\
  \midrule
  \endhead
  Rejestracja TRENERA z poprawnymi danymi            & newAccount & 200 \\
  Rejestracja KLIENTA z istniejącym trenerem         & newAccount & 200 \\
  Rejestracja KLIENTA bez trenera                    & newAccount & 200 \\
  Niepoprawny format adresu e-mail                   & newAccount & 400 \\
  Hasło krótsze niż 6 znaków                         & newAccount & 400 \\
  Zduplikowany adres e-mail                          & newAccount & 409 \\
  Nieprawidłowa rola (np.\ ADMIN)                    & newAccount & 400 \\
  Klient podaje ID nieistniejącego trenera            & newAccount & 400 \\
  Klient podaje ID użytkownika bez roli TRENER       & newAccount & 400 \\
  Logowanie z poprawnymi danymi                      & login      & 201 \\
  Logowanie ze złym hasłem                           & login      & 403 \\
  Logowanie z nieznanym adresem e-mail               & login      & 400 \\
  Logowanie bez pola \texttt{email}                  & login      & 400 \\
  Trener z przypisanymi klientami                    & clients    & 200 \\
  Trener bez klientów --- pusta lista                & clients    & 200 \\
  \bottomrule
\end{longtable}
\end{center}

Odpowiedź logowania weryfikowana jest dodatkowo pod kątem obecności pól \texttt{id},
\texttt{email}, \texttt{role} oraz \textbf{braku} pola \texttt{password} (sprawdzenie
działania \texttt{@JsonProperty(WRITE\_ONLY)}).

\subsection{Pokryte scenariusze --- \texttt{ExercisesControllerTest}}

Klasa testuje endpointy \texttt{GET /api/exercises}, \texttt{POST /api/history}
oraz \texttt{GET /users/\{id\}/history}.

\begin{center}
\begin{tabular}{p{7.5cm}cc}
  \toprule
  \textbf{Scenariusz} & \textbf{Endpoint} & \textbf{Status} \\
  \midrule
  Pobranie listy ćwiczeń (dane istnieją)          & GET /api/exercises      & 200 \\
  Pobranie ćwiczeń gdy baza pusta                 & GET /api/exercises      & 200 \\
  Dodanie wpisu historii z poprawnym body         & POST /api/history       & 201 \\
  Pobranie historii ćwiczeń użytkownika           & GET /users/\{id\}/history & 200 \\
  \bottomrule
\end{tabular}
\end{center}

\subsection{Wyniki}

\begin{center}
\begin{tabular}{lcc}
  \toprule
  \textbf{Klasa testowa} & \textbf{Liczba testów} & \textbf{Wynik} \\
  \midrule
  \texttt{UserControllerTest}      & 15 & \textcolor{green!60!black}{\textbf{PASS}} \\
  \texttt{ExercisesControllerTest} &  4 & \textcolor{green!60!black}{\textbf{PASS}} \\
  \midrule
  \textbf{Łącznie} & \textbf{19} & \textbf{0 błędów, 0 pominięć} \\
  \bottomrule
\end{tabular}
\end{center}

% ============================================================
\section{Bezpieczeństwo}

\subsection{Zastosowane mechanizmy}

\begin{itemize}
  \item \textbf{Walidacja danych wejściowych} --- backend weryfikuje poprawność formatu adresu e-mail (Apache Commons Validator) oraz minimalną długość hasła (wyrażenie regularne \texttt{\textbackslash w\{6,\}}).
  \item \textbf{Ochrona hasła w odpowiedziach API} --- pole \texttt{password} w encji \texttt{Users} oznaczone adnotacją \texttt{@JsonProperty(WRITE\_ONLY)}, co uniemożliwia jego zwrot w odpowiedziach GET.
  \item \textbf{Walidacja ról i relacji} --- backend sprawdza poprawność roli użytkownika i istnienie przypisywanego trenera przed zapisem.
\end{itemize}

\subsection{Znane ograniczenia bezpieczeństwa}
\label{sec:ograniczenia}

Poniższe braki zostały zidentyfikowane i są zaplanowane do eliminacji w kolejnych wersjach:

\begin{center}
\begin{tabular}{clc}
  \toprule
  \textbf{Lp.} & \textbf{Opis} & \textbf{Priorytet} \\
  \midrule
  1 & Hasła przechowywane w bazie jako plaintext (brak BCrypt) & Wysoki \\
  2 & Brak mechanizmu sesji / tokenów JWT & Wysoki \\
  3 & Komunikacja HTTP zamiast HTTPS & Wysoki \\
  4 & Brak blokady konta po wielokrotnych błędnych logowaniach & Średni \\
  5 & Brak wygasania sesji & Średni \\
  6 & Historia treningów nie synchronizowana z backendem & Niski \\
  \bottomrule
\end{tabular}
\end{center}

\textbf{Uwaga:} System w obecnej postaci przeznaczony jest wyłącznie do użytku w środowisku deweloperskim i akademickim. Przed wdrożeniem produkcyjnym wymagane jest usunięcie wszystkich wyżej wymienionych ograniczeń.

% ============================================================
\section{Polityka rozwoju oprogramowania}

System oparty jest na technologiach aktywnie wspieranych przez producentów:
\begin{itemize}
  \item Flutter / Dart --- Google; wsparcie LTS dla kanałów stable,
  \item Spring Boot 4.x --- VMware/Broadcom; wsparcie do co najmniej 2027,
  \item MySQL 8.0 --- Oracle; wsparcie do 2026 (Extended Support do 2030),
  \item Java 21 LTS --- wsparcie do 2029 (OpenJDK / Eclipse Temurin).
\end{itemize}

\textbf{Funkcje planowane w kolejnych wersjach:} szyfrowanie haseł (BCrypt), autentykacja JWT, HTTPS, synchronizacja historii Hive $\leftrightarrow$ backend, przegląd wideo sesji, moduł wiadomości trener--pacjent, obsługa kolejnych typów ćwiczeń w silniku detekcji pozy.

% ============================================================
\section{Spis dokumentów zewnętrznych}

\begin{center}
\begin{longtable}{lp{6cm}l}
  \toprule
  \textbf{Dokument} & \textbf{Opis} & \textbf{Źródło} \\
  \midrule
  \endfirsthead
  \toprule
  \textbf{Dokument} & \textbf{Opis} & \textbf{Źródło} \\
  \midrule
  \endhead
  Spring Boot Docs       & Dokumentacja frameworka Spring Boot 4.x & docs.spring.io \\
  Flutter Docs           & Dokumentacja SDK Flutter 3.x & docs.flutter.dev \\
  ML Kit Pose Detection  & Dokumentacja Google ML Kit & developers.google.com \\
  ExerciseDB API         & Dokumentacja zewnętrznego API ćwiczeń & rapidapi.com \\
  Hive Docs              & Dokumentacja bazy Hive dla Flutter & docs.hivedb.dev \\
  MySQL 8.0 Docs         & Dokumentacja MySQL Server 8.0 & dev.mysql.com \\
  Docker Docs            & Dokumentacja Docker Engine & docs.docker.com \\
  \bottomrule
\end{longtable}
\end{center}

% ============================================================
\appendix

\section{Pełna struktura katalogów projektu}

\dirtree{%
  .1 rehab\_app/.
  .2 frontend/.
  .3 lib/.
  .4 config/.
  .4 core/.
  .4 data/.
  .5 models/.
  .5 services/.
  .4 ui/.
  .5 screens/.
  .6 coach/.
  .7 coach\_main\_screen.dart.
  .7 patient\_list\_screen.dart.
  .7 patient\_detail\_screen.dart.
  .7 plan\_editor\_screen.dart.
  .7 video\_review\_screen.dart.
  .6 camera\_screen.dart.
  .6 catalog\_screen.dart.
  .6 exercise\_detail\_screen.dart.
  .6 login\_screen.dart.
  .6 session\_summary\_screen.dart.
  .6 workout\_history\_screen.dart.
  .5 widgets/.
  .4 vision/.
  .3 pubspec.yaml.
  .2 backend/.
  .3 src/main/java/com/example/demo/.
  .4 controllers/.
  .4 entities/.
  .4 repositories/.
  .4 DataInitializer.java.
}

\section{Kody odpowiedzi HTTP}

\begin{center}
\begin{tabular}{ll}
  \toprule
  \textbf{Kod} & \textbf{Znaczenie w kontekście API} \\
  \midrule
  200 OK          & Operacja zakończona sukcesem \\
  201 Created     & Zasób utworzony (logowanie, zapis historii) \\
  400 Bad Request & Nieprawidłowe dane wejściowe \\
  403 Forbidden   & Nieprawidłowe hasło \\
  409 Conflict    & Konflikt (e-mail już zajęty) \\
  \bottomrule
\end{tabular}
\end{center}

\end{document}